% review of existing alternatives
\section{REVIEW OF EXISTING ALTERNATIVES}

Currently, International School Manila is handling sponsorship billing complexity manually, each time the charge is entered; by looking up the caps and effective dates on spreadsheets; and by transferring files between independent platforms. This practice meets the requirement for posting in the current day and age but it also decentralizes authority, adds variability to the Bill-To decision and reduces the audit trail that downstream parties require. The implications are reduced processing speed, periodic reversals/reclassifications and more work to reconcile during the closing process.

In the “across the sector” there are three solution paths. The first approach is for institutions to adopt monolithic enterprise resource planning modules that integrate the functions of receivables, third-party billing and statements into a single suite, which work well when a charge finally reaches the receivables layer, but typically don't consider coverage to be computed in the charging interface \citep{EllucianBannerARGuide2025,PeopleSoftSFThirdParty,BannerARWorkbook}. Second, there are sponsor/third-party portals deployed in schools which offer statements, payments and sponsor self-service \citep{TouchNetSponsorPoint,NelnetSponsorBilling,UTKSponsorPoint,SponsorPaymentsPortal} - these tools have helped enhance the downstream experience, but they get whatever classification was done earlier at an entry point. Third, organisations strengthen existing processes by using templates, integrations and/or payments gateways, while now having processes that are more timely, but Letter of Guarantee interpretation is still spread across spreadsheets and scripts \citep{PayMyTuitionSponsored,PayMyTuitionNewswire}. The limits in suites for sponsor receipting and invoicing are similar, since they focus on billing instead of on the line level coverage logic at entry \citep{TribalEbsSponsor,TribalEbsManagePayments}. The receivables and Bill-To controls are still viable options for NetSuite, but it's not really built for evaluating LoG rules in the charging portal \citep{NetSuiteThirdParty}. For the Philippine context, the Ease of Paying Taxes framework also reinforces the need for proper documentation \citep{BIR_EOPT} in the form of Bill-To, a document that is auditable, whether it is for Internal Revenue or other tax authorities.

The deployed pilot system includes the following differences: Moving the coverage determination to the exact point where the charge is entered, and creating a Sponsor Master that is canonical and aligned with other systems of record. A deterministic rules engine analyzes each line in the Student Charging Portal and returns an explainable outcome (Covered, Split or Not Covered) that is then passed forward as the correct Bill-To and allocation to the ERP and then to the statement layer. This design helps to maintain existing platforms and avoids the use of ad hoc interpretation with a consistent and auditable logic.

\setlength{\LTleft}{0pt}
\setlength{\LTright}{\fill}
Table~\ref{table:existing-alternatives-summary} compares the proposed system with common sponsor billing, receivables, and payment alternatives so that the system gap can be assessed before the proposed design is discussed.

\begin{longtable}{>{\raggedright\arraybackslash}p{0.18\linewidth}
>{\raggedright\arraybackslash}p{0.14\linewidth}
>{\raggedright\arraybackslash}p{0.20\linewidth}
>{\raggedright\arraybackslash}p{0.22\linewidth}
>{\raggedright\arraybackslash}p{0.18\linewidth}}
\caption{Summary of existing alternatives}\label{table:existing-alternatives-summary}\\
\toprule
\textbf{System (License/Type)} & \textbf{Platform} & \textbf{Brief Description} & \textbf{Basic Functionalities} & \textbf{Relation to the Proposed System} \\
\midrule
\endfirsthead
\caption[]{Summary of existing alternatives (continued)}\\
\toprule
\textbf{System (License/Type)} & \textbf{Platform} & \textbf{Brief Description} & \textbf{Basic Functionalities} & \textbf{Relation to the Proposed System} \\
\midrule
\endhead
\midrule
\multicolumn{5}{r}{\textit{Continued on next page}}\\
\endfoot
\bottomrule
\endlastfoot
\textbf{Proposed System (ISM)} & .NET responsive web app; MSSQL; APIs & Coverage at entry plus canonical Sponsor Master & LoG rules in SCP; instant Covered/Split/Not Covered with reason; correct Bill-To and sponsor-parent allocations to ERP; clear statements; full audit trail & -- \\
\textbf{TouchNet SponsorPoint}\tnote{d} & Hosted portal & Sponsor/third-party billing portal for statements, payments, reconciliation & Sponsor dashboards; electronic statements; online payments; ERP sync & Downstream billing/sponsor access; assumes coverage decided upstream; no LoG rules at SCP entry \\
\textbf{Nelnet Sponsor Billing \& Payments}\tnote{e} & Hosted portal & Sponsor portal with configurable billing cycles and ERP sync & Billing cycles; statements; payments; account management & Improves billing/payments; inherits upstream classification; not a coverage engine at entry \\
PayMyTuition Sponsored Payments\tnote{f} & Payments gateway/module & Sponsored payments with automated billing and reconciliation & Automated sponsor billing; reconciliation; real-time integrations & Payments/reconciliation focus; does not compute coverage at charge entry \\
Ellucian Banner A/R (Third-Party)\tnote{a,c} & ERP module & Receivables module supporting third-party contracts and billing runs & Detail codes; contracts; invoice runs; collections; audit controls & Centralizes receivables; expects eligibility/coverage decisions pre-posting \\
\end{longtable}

\subsection{Summary of Existing Alternatives}

While there are existing alternatives that enhance posting, invoicing, payments and collections, they cannot address the key driver of the inconsistencies, which is the manual and context-specific interpretation of \textbf{Letters of Guarantee} within the charge interface. There is fragmentation and variability in \textbf{Bill-To} decisions, and audit trails - across ERP modules, sponsor portals, payment gateways and student-suite extensions, coverage is assumed to be settled upstream, or inherited from manual entry. Thus, implementing any of these systems “as is” would perpetuate ongoing occurences of reversals, re-classifications, and mismatches in real-time coverage and the charge entry of a single authoritative Bill-To in NetSuite and the Online Billing System.

The deployed pilot system works to overcome this limitation by moving the decision locus to charge entry and by implementing a canonical Sponsor Master and synchronizing this system with \textbf{PowerSchool}, \textbf{Student Charging Portal}, \textbf{NetSuite}, and \textbf{Online Billing System}. Each line is compared to codified provisions, items/categories, caps, effective dates and per student overrides by a deterministic rules engine and an outcome with a reason code is immediately returned before the charge is saved. The first determination is the final Bill-To/student-sponsor determination in \textbf{NetSuite} and is passed to statements forward without further reclassification. Rules, actors, time-stamps and context are captured with each evaluation to create a full line-level audit chain aligned with Ease of Paying Taxes requirements.

The unique combination of a single source of truth for sponsors and real-time, explainable coverage logic provides three unique benefits: a single decision recorded at line-level, authoritative propagation of \textbf{Bill-To} and allocations across ledgers and statements, and better audit readiness without replatforming. This synergy of decision locus, data governance and integration provides a viable path to greater reliability, faster processing and continuous compliance in the context of International School Manila.
